\section*{Bab \ref{ch:b3:plant-molecular-physiology} --- Fisiologi Tumbuhan Molekuler dan Tanggapan Tekanan}

\begin{solution}{exo:b3:plant-molecular-physiology:1}
\omterm{def:b3:plant-molecular-physiology:photoreceptors}{Fitokrom}, merah \qty{660}{nm} dan merah-jauh \qty{730}{nm}:
pengurangan \omterm{def:b3:plant-molecular-physiology:photoreceptors}{etiolasi}, \omterm{prop:b3:plant-molecular-physiology:photoequilibrium}{penghindaran naungan}, perkecambahan. \omterm{def:b3:plant-molecular-physiology:photoreceptors}{Kriptokrom}, biru
\qty{450}{nm}: pengurangan \omterm{def:b3:plant-molecular-physiology:photoreceptors}{etiolasi}, jamnya, pembungaan. \omterm{def:b3:plant-molecular-physiology:photoreceptors}{Fototropin}, biru:
pembengkokan ke arah cahaya, membukanya stomata, gerakan kloroplas.
Klorofilnya mengukur cahaya sebagai tenaga lalu tidak melakukan apa pun dengan
keterangannya; fotoreseptornya membaca warna, arah, dan lamanya pada
fluks sejuta kali lebih rendah lalu mengubah pengungkapan gennya --- yakni
perbedaan antara memakan cahayanya dan membacanya.
\end{solution}

\begin{solution}{exo:b3:plant-molecular-physiology:2}
Tiap \omterm{def:b3:endocrinology:hormone}{hormonnya} mendorong pengubikuitinasian dan pemusnahan satu
penekan: \omterm{def:b3:plant-molecular-physiology:hormones}{auksin} merekatkan \omterm{def:b3:plant-molecular-physiology:hormones}{TIR1} pada protein Aux/IAA, \omterm{def:b3:plant-molecular-physiology:hormones}{giberelin} yang terikat
pada GID1 menghantarkan protein \omterm{def:b3:plant-molecular-physiology:hormones}{DELLA} ke sebuah ligase F-box, jasmonat merekatkan
COI1 pada protein JAZ. Penekannya sudah duduk pada gen sasarannya
dengan pengaktifnya bersiaga di sebelahnya, sehingga tidak perlu ada protein baru yang
dibuat: memusnahkan penekannya, yang dikerjakan proteasomnya dalam
hitungan menit, menyalakan gennya seketika.
\end{solution}

\begin{solution}{exo:b3:plant-molecular-physiology:3}
ABA mengikat PYR/PYL; kompleksnya menghambat fosfatase PP2C-nya;
kinase SnRK2-nya (OST1), yang tidak lagi didefosforilasi, menjadi aktif lalu
memfosforilasi saluran anion SLAC1-nya; anionnya pergi, selaputnya
terdepolarisasi, saluran kalium keluarnya membuka, K$^{+}$ lalu air
pergi, turgornya turun dan \omterm{prop:b3:plant-molecular-physiology:stomata}{sel penjaganya} melemas bersama-sama. Tiap mata rantainya dapat
diputus: tanpa reseptor, sebuah fosfatase yang tidak dapat dihambat (yakni
mutan \emph{abi1} yang klasik), tanpa OST1 atau tanpa SLAC1 --- masing-masing memberi sebuah
tumbuhan layu yang stomatanya tetap terbuka pada kekeringan.
\end{solution}

\begin{solution}{exo:b3:plant-molecular-physiology:4}
PTI mengenali molekul yang lazim pada seluruh kelas mikroba
(flagelin, kitin) dengan kinase reseptor di permukaan selnya lalu berakhir
pada sebuah dinding yang diperkuat, stomata yang tertutup, pengungkapan gen pertahanan, dan
pertumbuhan yang lebih lambat, biasanya tanpa kematian sel. ETI mengenali sebuah
protein efektor tertentu atau kerusakannya dengan sebuah \omterm{def:b3:plant-molecular-physiology:immunity}{reseptor NLR} di dalam sel lalu
berakhir pada kematian hipersensitif sel yang terinfeksi dan sebuah tanda bahaya
menyeluruh.
\end{solution}

\begin{solution}{exo:b3:plant-molecular-physiology:5}
Cahaya siang $0.87\times 1.15/1.75 = 0.57$; sehelai daun $0.87\times 0.4/1.0 =
0.35$; kanopi rapat $0.87\times 0.1/0.7 = 0.12$; lampu pijar
$0.87\times 0.7/1.3 = 0.47$. Sebuah lampu berfilamen kaya merah-jauh, sehingga
$\varphi$-nya di bawah nilai cahaya siangnya dan semainya membaca naungan sebagian:
ia memanjang. Lampu putih sejuk dan LED tanpa merah-jauh melakukan
sebaliknya.
\end{solution}

\begin{solution}{exo:b3:plant-molecular-physiology:6}
Dinding $5.0$, sitosol $7.5$: $(1 + 10^{2.75})/(1 + 10^{0.25}) = 563/2.78
= 203$. Dinding $4.5$: $563/(1 + 10^{-0.25}) = 563/1.56 = 360$ ---
mengasamkan dindingnya menaikkan pecahan netral di luarnya dan pengambilannya.
Sitosol $6.5$ dengan dinding $5.5$: $(1 + 10^{1.75})/6.62 = 57/6.62 = 8.6$:
penjebakannya turun lima kali lipat, \omterm{def:b3:plant-molecular-physiology:hormones}{auksinnya} bocor keluar dari tiap mukanya, dan
landaian kutubnya mendatar --- salah satu sebab akar tanpa oksigen kehilangan arahnya.
\end{solution}

\begin{solution}{exo:b3:plant-molecular-physiology:7}
\qty{1}{cm/h} $= \qty{10000}{\micro m/h}$: yakni $200$ sel tiap jam,
\qty{18}{s} di dalam masing-masing. Pembauran melintasi \qty{50}{\micro m}: $t =
(5\times 10^{-5})^{2}/(2\times 5\times 10^{-10}) = \qty{2.5}{s}$. Bagian
dalam selnya dilintasi dalam hitungan detik; kebanyakan dari \qty{18}{s} itu
dihabiskan menunggu untuk diusung keluar lewat PIN pada selaput
basalnya, yang menjadi langkah pembatas lajunya.
\end{solution}

\begin{solution}{exo:b3:plant-molecular-physiology:8}
$E/A = 1.6\,\Delta w/\Delta c = 1.6\times\num{16000}/150 = 171$ molekul
air tiap CO$_{2}$. C$_{4}$, $\Delta c = 250$: yakni $102$. Hari kering,
$\Delta w = 24$: yakni $256$.
\end{solution}

\begin{solution}{exo:b3:plant-molecular-physiology:9}
Pembekuan mengeringkan selnya dengan menarik airnya ke dalam es luar selnya, sehingga
dingin dan kekeringan keduanya tekanan pengeringan. Keduanya berbagi pesuruh
kedua (yakni ABA, lonjakan kalsium, dan oksigen reaktif), faktor
transkripsi (keluarga CBF/DREB dipicu keduanya), dan produknya:
dehidrin dan protein LEA yang melindungi selaput dan protein dari
kehilangan air, zat terlarut serasi (prolin, gula) yang menahan airnya, serta
antioksidan. Sebuah tumbuhan yang telah membuatnya bagi satu tekanan punya semuanya bagi
tekanan yang lain.
\end{solution}

\begin{solution}{exo:b3:plant-molecular-physiology:10}
Pada \qty{1}{\%} matahari penuh, $k_{1} + k_{2} = \qty{0.005}{s^{-1}}$: yakni 95~\%
kesetimbangannya sesudah $3/k = \qty{600}{s}$, yakni sepuluh menit, terhadap enam
detik pada tengah hari. Delapan jam adalah empat paruh hidup: $1/16 \approx
\qty{6}{\%}$ Pfr-nya tersisa. Sebuah kilatan singkat menetapkan $\varphi$ tetapi
Pfr-nya lalu meluruh, sedangkan sehari menjaganya tetap tinggi berjam-jam; tanggapan
yang menuntut Pfr bekerja berjam-jam (perkecambahan dan pengurangan \omterm{def:b3:plant-molecular-physiology:photoreceptors}{etiolasi})
mengintegralkan pemaparannya, sehingga sebuah biji yang tersingkap sebuah bajak sesaat tidak menanggapinya
seperti biji yang tergeletak di permukaannya.
\end{solution}

\begin{solution}{exo:b3:plant-molecular-physiology:11}
Sepuluh tapak tiap hari pada $100$ sel masing-masing: yakni \num{1000} sel, $10^{-4}$
daunnya; sepanjang semusim sepersekian persen, terhadap kehilangan
daunnya bila satu infeksi menyebar tanpa dikekang. Sel yang mati berongkos nyaris
tidak ada lalu membawa makanan patogennya bersamanya --- bagi sebuah biotrof,
yang memerlukan sel yang hidup. Sebuah nekrotrof makan dari jaringan mati:
\omterm{def:b3:plant-molecular-physiology:immunity}{tanggapan hipersensitifnya} menyerahkan sesantapan dan sebuah pijakan kepadanya, dan sebagian
(\emph{Botrytis}, \emph{Sclerotinia}) menyekresikan racun yang memancingnya
dengan sengaja; terhadapnya tumbuhannya bergantung pada pertahanan yang diperantarai jasmonat
dan bukan pada kematian sel.
\end{solution}

\begin{solution}{exo:b3:plant-molecular-physiology:12}
Produksinya pada laju $p$ dari jam 12 sampai 16. Hari panjang, tersinari sepanjang waktu,
$k = \ln 2/4 = \qty{0.17}{h^{-1}}$: arasnya pada 16 adalah $(p/k)(1 -
\mathrm{e}^{-4k}) = (p/0.17)(0.5) = 2.9p$. Hari pendek, gelap sejak jam
10, $k = \qty{1.39}{h^{-1}}$: $(p/1.39)(1 - \mathrm{e}^{-5.5})
\approx 0.72p$. Empat kali lebih banyak CONSTANS pada hari panjangnya. Sebuah ambang
di antaranya --- katakanlah $2p$, yang diperlukan untuk mengaktifkan \emph{FT} --- dilintasi hanya
ketika jendela produksi jamnya berbarengan dengan cahayanya: kebetulan
sebuah irama dalam dengan hari luarnya mengubah sebuah panjang hari yang berjenjang
menjadi sebuah keputusan berbunga yang sama sekali ada atau tidak ada.
\end{solution}

\begin{solution}{pb:b3:plant-molecular-physiology:1}
\textbf{1.} Cahaya siang $0.87\times 1.15/1.75 = 0.57$; kanopi $0.87\times
0.2/0.8 = 0.22$.
\textbf{2.} $r = 2(1 - \varphi)$: \qty{0.86}{mm/h} pada cahaya siang,
\qty{1.57}{mm/h} di dalam naungan; sepanjang \qty{48}{h}: $41$ terhadap \qty{75}{mm},
yakni \qty{34}{mm} tambahan.
\textbf{3.} Merahnya turun menjadi $0.1\times 1.15 = 0.115$ dari merah-jauhnya, yang
turun menjadi $0.8$: $\zeta = 0.144$, $\varphi = 0.87\times 0.144/0.744 =
0.17$. Ya: sehelai daun tunggal menjatuhkan $\varphi$ dari $0.57$ ke $0.17$, jauh di
dalam jangkauan \omterm{prop:b3:plant-molecular-physiology:photoequilibrium}{penghindaran naungannya}.
\textbf{4.} Kedua laju perubahan fotonya berskala dengan fluksnya, sehingga
nisbahnya tidak. Di bawah awan, spektrumnya nyaris tak berubah dan
$\varphi$-nya bertahan dekat $0.57$: sebuah semai di tempat terbuka pada hari yang muram tidak
memanjang --- ia membaca tetangganya, bukan kecerahannya.
\textbf{5.} \qty{95}{\%} sesudah $3/(k_{1} + k_{2})$: yakni \qty{6}{s} pada matahari
penuh, \qty{600}{s} pada \qty{1}{\%}.
\textbf{6.} \qty{8}{h} $=$ empat paruh hidup, $1/16 = \qty{6.3}{\%}$;
\qty{14}{h} $=$ tujuh, $1/128 = \qty{0.8}{\%}$. Pfr sisanya saat fajar
dapat melaporkan panjang malamnya, tetapi pembalikannya adalah sebuah reaksi termal
yang berjalan lebih cepat pada malam yang hangat, dan nilai awalnya bergantung pada
cahaya harinya: tumbuhan mengukur panjang malamnya dengan jamnya sebagai gantinya.
\textbf{7.} Dinding: $1/(1 + 10^{0.75}) = 0.15$; sitosol: $1/(1 +
10^{2.45}) = 0.0035$.
\textbf{8.} $(1 + 282)/(1 + 5.6) = 283/6.6 = 43$. Sitosolnya
$43\times 0.1 = \qty{4.3}{\micro M}$.
\textbf{9.} $\qty{1}{cm/h}/\qty{100}{\micro m} = 100$ sel tiap jam,
\qty{36}{s} masing-masing.
\textbf{10.} \qty{5}{h}. Pembengkokannya mulai dalam sejam sebab
pengagihan ulang yang berarti adalah pengagihan ke samping, melintasi semilimeter pucuknya,
dan tanggapannya setempat; aliran jarak jauhnya menetapkan pasokannya,
bukan ketepatan waktunya.
\textbf{11.} Sayap bawahnya: $60/50 = 1.2$, yakni \omterm{def:b3:plant-molecular-physiology:hormones}{auksin} \qty{20}{\%} di atas normanya,
pemanjangannya $0.8$; sayap atasnya: $0.8$ dari normanya, pemanjangannya $1.2$. Sayap
atasnya melampaui sayap bawahnya: akarnya melengkung ke bawah.
\textbf{12.} Sel tunasnya terletak di bawah optimum \omterm{def:b3:plant-molecular-physiology:hormones}{auksinnya}, sehingga \omterm{def:b3:plant-molecular-physiology:hormones}{auksin}
tambahan pada sayap bawahnya mendorong pertumbuhannya: sayap bawahnya melampaui
sayap atasnya dan tunasnya melengkung ke atas.
\textbf{13.} $E = 0.3\times 16 = \qty{4.8}{mmol\,m^{-2}\,s^{-1}}$.
\textbf{14.} $A = (0.3/1.6)\times 150 = \qty{28}{\micro mol\,m^{-2}\,s^{-1}}$;
$E/A = 4800/28 = 171$.
\textbf{15.} Luasnya \qty{0.002}{m^2}, \num{43200}~s: airnya $4.8\times
10^{-3}\times 0.002\times\num{43200} = \qty{0.41}{mol} = \qty{7.5}{g}$;
CO$_{2}$-nya $28\times 10^{-6}\times 0.002\times\num{43200} =
\qty{2.4e-3}{mol}$, yakni $2.4\times 10^{-3}/6\times 180 =
\qty{0.073}{g}$ glukosa --- yakni seratus gram air bagi segram
gula.
\textbf{16.} $E = 0.05\times 16 = \qty{0.8}{mmol\,m^{-2}\,s^{-1}}$,
$A = \qty{4.7}{\micro mol\,m^{-2}\,s^{-1}}$, nisbahnya tetap $171$. Tumbuhannya
telah menghemat enam pertujuh airnya lalu menyerahkan enam pertujuh
karbonnya: menutup katupnya mengubah lajunya, bukan harganya.
\textbf{17.} C$_{4}$: $\Delta c = 250$, $E/A = 4800/46.9 = 102$. CAM pada
malam hari: $E = 0.3\times 5 = \qty{1.5}{mmol}$, $A = 28$: $E/A = 53$.
\textbf{18.} $E/A = 1.6\,\Delta w/\Delta c$: kehantarannya meniadakan diri
sebab kedua gasnya melewati liang yang sama. Tumbuhannya dapat mengubah
landaiannya --- memekatkan CO$_{2}$ di dalamnya (C$_{4}$), membuka ketika udaranya
lembap (CAM, fajar), mendinginkan daunnya, menebalkan lapisan batasnya dengan bulu
--- tetapi bukan fisika dua gas yang berbagi satu lubang.
\textbf{19.} \num{1000} sel tiap hari, yakni $10^{-4}$ daunnya.
\textbf{20.} \num{60000} sel dalam 60 hari, yakni \qty{0.6}{\%}. Satu lolosan tiap
hari yang memusnahkan \qty{5}{\%} akan memakan habis daunnya dalam dua puluh hari:
\omterm{def:b3:plant-molecular-physiology:immunity}{tanggapan hipersensitifnya} berongkos seperseratus ongkos sebuah infeksi
yang lolos.
\textbf{21.} $0.2\times 5 = \qty{1}{\%}$ fotosintesisnya. Bila dinyalakan
selalu ia akan berongkos \qty{5}{\%} dan tumbuhannya akan dilampaui
tetangga yang membelanjakannya pada daun; pertahanan terimbas berbalas hanya ketika
ancamannya hadir.
\textbf{22.} Ia memperoleh ketaktampakan bagi NLR itu lalu menginfeksi ragam yang
kebal; ia kehilangan apa pun yang dilakukan efektornya (menekan PTI), sehingga ia
lebih lemah pada tumbuhan tanpa NLR itu. Populasi tumbuhannya berakhir
memihak NLR terhadap efektor yang tersisa, dan gen kekebalan yang lama,
yang kini tak berguna, merosot: yakni pendauran bergantung-kekerapan gen kedua
pihaknya.
\textbf{23.} Tiruannya (koronatin) mengaktifkan jalur jasmonatnya,
yang melawan jalur salisilatnya; salisilat adalah pertahanan
terhadap biotrof seperti bakterinya sendiri, dan tiruannya juga
membuka kembali stomata yang telah ditutup tumbuhannya. Patogennya memutar satu lengan
sistem imunnya melawan lengan yang lain.
\textbf{24.} Sebuah gen kekebalan tunggal atas berjuta-juta tumbuhan yang serupa
adalah sebuah seleksi yang seragam: mutan mana pun yang kehilangan atau mengubah efektornya
menyebar ke seluruh tanamannya dalam beberapa musim. Pemulianya menumpuk beberapa
gen kekebalan sehingga beberapa efektornya harus hilang sekaligus, mencampur
ragamnya, menggilir gennya bertahun-tahun, dan memihak kekebalan terpicu pola
yang luas yang tidak terhindari satu mutasi pun.
\textbf{25.} $\varphi \approx 0.22$ di bawah kanopinya; \omterm{def:b3:plant-molecular-physiology:hormones}{auksin} dalam :
luar $\approx 43$; sekitar $170$ molekul air yang hilang tiap CO$_{2}$
yang difiksasi.
\end{solution}
